\documentclass[a4paper,12pt]{article}
\usepackage{../common/macros}
\usepackage{layout}
\usepackage{url}
\usepackage{booktabs}

\title{マルチメディアコミュニケーションシステム 課題1\\
{\normalsize 一般電話と相互接続可能なスマートフォン通話アプリの調査}}
\author{S2620154　吉田 将衛　システム知能クラスタ}

\renewcommand{\today}{\the\year 年\the\month 月\the\day 日}
\date{\today}

\begin{document}

\maketitle

% ===== Rakuten Link =====
\section{Rakuten Link（楽天モバイル）}\cite{rakuten-link-service}

\subsection{料金体系}

Rakuten Linkどうしでの通話，メッセージは無料である．
また，国内への通話に関しては，Rakuten Linkを利用していない相手への通話も無料である．
逆に，海外への通話に関しては，国や地域によっては従量課金制であるため，注意が必要である．
北米エリアは30秒で34円，韓国，台湾，中国，香港，マカオは30秒で57円，イギリス，イタリア，オーストリア，オランダ，グリーンランド，スイス，スペインは30秒で108円，アフリカ地域は30秒で180円，中東地域は30秒で148円，中南米は30秒で111円または148円である．
電話の着信については，国内外問わず無料である．

\subsection{通話品質}

ギガ無制限エリアでは，ギガ無制限で高速通信を利用可能である．
Wi-Fi接続時は，Wi-Fiの通信品質に依存するため，LTE接続での利用が推奨されている．
VPNやデュアルSIMでの利用時は楽天の動作保証対象外となる．
海外ローミングを使用した場合，データ利用料が2GBを超過すると，通信速度が最大128kbpsに制限されるため，通話品質が低下する可能性がある．

\subsection{メリット}

自身の電話番号から他の携帯電話会社の電話番号（固定電話を含む）に無料で通話できる．
海外の対象国・地域であれば日本国内へな無料で電話がかけられる．
メッセージの送受信が無料である．Rakuten Linkどうしであれば，メッセージはもちろん，写真や動画，ファイルも送受信でき，最大100人のグループメッセージも楽しめる．
楽天モバイル契約者限定のキャンペーンや特典を受け取ることができたり，スマホ決済サービスに簡単にアクセスできたり，楽天ポイントを貯めたり使ったりできる．

\subsection{デメリット・利用制限}

対象外の電話番号への通信は自動でOS標準の電話アプリに切り替わるため，有料となる．
以下が対象外の電話番号である．
\begin{itemize}
  \item 0570（ナビダイヤル）
  \item 171（災害伝言ダイヤル）
  \item 188（消費者ホットライン）
  \item \#7119（総務省消防庁）etc.
\end{itemize}
iOS版は，Rakuten Linkでは着信ができないため，iOSの標準の電話アプリでの着信となる．
Wi-Fi，VPN，デュアルSIMでの利用は楽天の動作保証対象外となる．
SMSの送信上限は一日あたり最大200通までである．
Webサイトのリンクから発信した場合や着信履歴からの折り返しは自動的にOS標準の電話アプリに切り替わるため，有料となる．
毎月2GB付与される海外ローミングエリアの高速データ容量が余った場合でも，高速データ容量を翌月に繰り越すことはできない．
エリア境目などで，実際には4Gサービスが提供されているにもかかわらず，「5G」と表示される場合がある．

% ===== G-Call 050 =====
\section{G-Call 050（ジー・コミュニケーション）}\cite{gcall050-service}

\subsection{料金体系}

月額の基本料金は550円（税込）である．
基本的には，G-Call050どうしであれば，国内外問わず無料で通話できる．
国内への通話料金に関しては，G-Callどうしであれば，固定電話への料金は3分あたり8.8円（税込）であり，携帯電話への通話料金は1分あたり17.6円（税込）である．
海外への通話料金は1分ごとの課金となる．
アメリカ，ハワイは5円，アルゼンチン，カナダ，メキシコは10円である．
中国は23円，韓国30円，香港35円，台湾10円である．
オーストラリアは10円，グアム40円，サイパン20円である．
なお，ヨーロッパやアフリカ地域は対応していない．
留守電機能に関しては無料で利用できる．

\subsection{通話品質}

インターネットを利用した回線であるため，ネットワーク環境が悪い場合は雑音が入ったり，音が途切れる場合がある．通話のセキュリティを完全に保証するものではない．

\subsection{メリット}

G-Call050を持てば，世界中のどこにいても世界中で受発信でき，G-Call050どうしであれば，無料である．
世界中どこにいても番号通知で発信できる．
音声SIM契約が不要で，データ専用SIMでも通話が可能である．
iPhone，Android，iPad，iPod touchなど，様々な端末で利用できる．
G-Callポイントが請求金額の1\%として貯まり，1ポイント1円で対象サービスの支払いに利用できる．
どのキャリア・MVNO（ドコモ・au・ソフトバンク・楽天等）でも利用可能である．
アプリがバックグラウンドや非起動時でもプッシュ通知で着信ができる．

\subsection{デメリット・利用制限}

ネット接続状態によって通話品質が悪くなる可能性がある．
日本の電話番号以外へ発信する場合は国と地域が制限される．
国際電話利用の際は番号通知は保証されない．
利用環境はCloud Softphoneアプリが動作する環境に限られる．また，Cloud Softphoneアプリのサポートはベストエフォートであるため，通知を行うサーバなどに障害があっても原因や対策の問い合わせはできない．
緊急通報（110番，118番，119番）及び3桁番号サービス（104/115/117/177等）への通話は，利用できない．
通話転送・保留・暗号化（SSL/TLS）等は非対応である．

% ===== 参考文献 =====
\bibliographystyle{jplain}
\bibliography{refs}

\end{document}
